\section{Post-Confirmatory Diagnostics}
\label{sec:diagnostics}

The analyses in this section were specified or conducted after the held-out
result. They explain and stress-test the frozen policy; they do not alter its
selection or rescue the failed absolute gate.

\subsection{Does the common-seed championship help?}

The championship addressed a training-only concern: each optimizer run had
selected its six finalists under different seeds. We therefore compared the
champion with the five policies that each run would have selected on its own,
using four new common seed blocks on the ten open method-v2 development
families. The complete design contains \MechanismGameCount{} games (six
methods, ten families, four seed blocks, and reciprocal slots).

The champion scores \MechanismChampionScore{}, compared with
\MechanismLocalAverageScore{} for the equal average of the
five run-local policies. The difference is \MechanismEstimate{}, but its
family-clustered 95\% interval
[\MechanismCILower{},~\MechanismCIUpper{}] includes zero. All five pairwise
point differences favor the champion (\MechanismPairwiseMinimum{} to
\MechanismPairwiseMaximum{}), while every pairwise
interval crosses zero. This evidence is consistent with common-seed selection
being useful, but ten families do not statistically resolve the mechanism.

\subsection{Which policy group matters?}

We formed five policies that each revert one group of champion settings to the
default: defense-radius growth, emergency defense, forced attack, scouting, or
the joint attack-composition/strategic-plan group. A second 480-game
common-seed panel compared all reverts with the champion on the same ten open
families. The primary champion-minus-equal-revert-average estimate is
\ComponentAverageEstimate{} with 95\% interval
[\ComponentAverageCILower{},~\ComponentAverageCIUpper{}].

Figure~\ref{fig:components} reports every individual contrast. Reverting
infantry+rush to assault/macro gives by far the largest decline:
\StrategyEstimate{}, with ordinary 95\% interval
[\StrategyOrdinaryLower{},~\StrategyOrdinaryUpper{}]. However, its
prespecified five-comparison Bonferroni familywise interval is
[\StrategyFamilywiseLower{},~\StrategyFamilywiseUpper{}], narrowly including
zero. The defense-growth and forced-attack reverts have slightly higher point
scores than the champion, emergency defense is near zero, and the scouting
revert is endpoint-identical in all 80 paired games. Thus the joint strategy
group is the dominant observed signal, but no individual component clears the
multiplicity-controlled criterion.

\begin{figure}[t]
  \centering
  
\begin{tikzpicture}
\begin{axis}[
  width=0.90\linewidth,
  height=5.2cm,
  xmin=-0.15, xmax=0.72,
  ymin=0.4, ymax=5.6,
  xtick={-0.1,0,0.1,0.3,0.5,0.7},
  ytick={1,2,3,4,5},
  yticklabels={{Defense growth},{Emergency defense},{Forced attack},{Scouting},{Infantry + rush}},
  xlabel={Champion $-$ one-group revert score},
  axis line style={draw=ChronoInk},
  tick label style={text=ChronoInk,font=\scriptsize},
  label style={text=ChronoInk,font=\small},
  grid=major,
  grid style={draw=ChronoGrid},
  clip=false
]
\addplot[draw=ChronoInk,dashed,line width=0.7pt] coordinates {(0,0.4) (0,5.6)};
\draw[ChronoBlue,line width=1.1pt] (axis cs:-0.08997983,1.12) -- (axis cs:0.01497983,1.12); \addplot[only marks,mark=*,mark size=1.8pt,draw=ChronoBlue,fill=ChronoBlue] coordinates {(-0.03750000,1.12)}; \draw[ChronoOrange,line width=1.1pt] (axis cs:-0.11289299,0.88) -- (axis cs:0.03789299,0.88); \addplot[only marks,mark=square*,mark size=1.7pt,draw=ChronoOrange,fill=ChronoOrange] coordinates {(-0.03750000,0.88)}; \draw[ChronoBlue,line width=1.1pt] (axis cs:-0.04727743,2.12) -- (axis cs:0.05977743,2.12); \addplot[only marks,mark=*,mark size=1.8pt,draw=ChronoBlue,fill=ChronoBlue] coordinates {(0.00625000,2.12)}; \draw[ChronoOrange,line width=1.1pt] (axis cs:-0.07064799,1.88) -- (axis cs:0.08314799,1.88); \addplot[only marks,mark=square*,mark size=1.7pt,draw=ChronoOrange,fill=ChronoOrange] coordinates {(0.00625000,1.88)}; \draw[ChronoBlue,line width=1.1pt] (axis cs:-0.05867609,3.12) -- (axis cs:0.03367609,3.12); \addplot[only marks,mark=*,mark size=1.8pt,draw=ChronoBlue,fill=ChronoBlue] coordinates {(-0.01250000,3.12)}; \draw[ChronoOrange,line width=1.1pt] (axis cs:-0.07883699,2.88) -- (axis cs:0.05383699,2.88); \addplot[only marks,mark=square*,mark size=1.7pt,draw=ChronoOrange,fill=ChronoOrange] coordinates {(-0.01250000,2.88)}; \addplot[only marks,mark=diamond*,mark size=2pt,draw=ChronoMuted,fill=ChronoMuted] coordinates {(0.00000000,4)}; \draw[ChronoBlue,line width=1.1pt] (axis cs:0.09556151,5.12) -- (axis cs:0.56693849,5.12); \addplot[only marks,mark=*,mark size=1.8pt,draw=ChronoBlue,fill=ChronoBlue] coordinates {(0.33125000,5.12)}; \draw[ChronoOrange,line width=1.1pt] (axis cs:-0.00734223,4.88) -- (axis cs:0.66984223,4.88); \addplot[only marks,mark=square*,mark size=1.7pt,draw=ChronoOrange,fill=ChronoOrange] coordinates {(0.33125000,4.88)};
\end{axis}
\end{tikzpicture}

  \caption{Champion-minus-one-group-revert effects. Blue circles/lines are
  ordinary family-clustered 95\% intervals; orange squares/lines are
  Bonferroni familywise-95\% intervals. The scouting contrast (gray diamond)
  has zero endpoint difference and zero observed variance, so no artificial
  interval is drawn.}
  \label{fig:components}
\end{figure}

\subsection{Terminal army and economy pattern}

A deterministic analyzer validated and reduced all \TerminalRawRecordCount{} raw completion
records from the confirmatory and two diagnostic panels. Across confirmatory
pairs, the champion improves its candidate-minus-opponent terminal balance by
\ConfirmatoryTerminalCombatantDifference{} combatants and
\ConfirmatoryTerminalUnitDifference{} total units while holding
\ConfirmatoryTerminalCreditReduction{} fewer credits.
Outcome changes mechanically alter terminal snapshots, so we also examine the
\StableDrawPairCount{} pairs that remain tick-cap draws under both methods.
They have identical score and duration, yet the champion ends with
\StableDrawCombatantDifference{} more relative combatants,
\StableDrawUnitDifference{} more relative units, and
\StableDrawCreditDifference{} relative credits.

The strategy-revert panel has the same qualitative pattern: the champion has
\StrategyTerminalCombatantDifference{} more relative combatants,
\StrategyTerminalUnitDifference{} more relative units, and
\StrategyTerminalCreditReduction{} fewer credits overall. These endpoints
are consistent with converting banked
resources into fielded combat power. They do not establish when or how that
conversion occurred, because the preserved games contain no action trace,
production timeline, combat log, or intermediate state trajectory. Likewise,
endpoint identity for the scouting revert does not prove identical internal
actions.
